The GLWE cryptosystem is a generalized form to encompass both the LWE and RLWE cryptosystems. The GLWE cryptosystem's ciphertext is a tuple $(\{A_i\}_{i=0}^{k-1}, B)$, where $B = \sum\limits_{i=0}^{k-1}{(A_i \cdot S_i)} + \Delta \cdot M + E$. The public mask $\{A_i\}_{i=0}^{k-1}$ and the secret key $\{S_i\}_{i=0}^{k-1}$ are a list of $k$ $(n-1)$-degree polynomials, each. The message $M$ and the noise $E$ are $(n-1)$-degree polynomials, each. Like in LWE and RLWE, a random public mask $A$ is created for each ciphertext, whereas the same secret key $S$ is used for all ciphertexts. In this section, we denote each ciphertext instance as $(\{A_i\}_{i=0}^{k-1}, B)$ instead of $(\{A_i\}_{i=0}^{k-1, \langle j \rangle}, B^{\langle j \rangle})$ for simplicity.




\subsection{Setup}
Let $t$ be the size of plaintext, and $q$ the size of ciphertext, where $t < q$ ($t$ is much smaller than $q$) and $t | q$ (i.e., $t$ divides $q$). Randomly pick a list of $k$ $(n-1)$-degree polynomials as a secret key, where each polynomial coefficient is a randomly picked ternary number in $\{-1, 0, 1\}$ (i.e., $\{S_i\}_{i=0}^{k-1} \xleftarrow{\$} \mathcal{R}_{\langle n, \textit{tern} \rangle}^k$). Let $\Delta = \left\lfloor\dfrac{q}{t}\right\rfloor$ be the scaling factor of plaintext.

Notice that GLWE's setup parameters are similar to that of RLWE. One difference is that $S$ is not an $(n-1)$-degree polynomial encoding $n$ secret coefficients, but a list of $k$ such $(n-1)$-degree polynomials encoding total $n \cdot k$ secret coefficients.


\subsection{Encryption}
\label{subsec:glwe-enc}

Polynomial representatives have degree at most $n-1$, including when leading coefficients vanish. Noise coefficients are integer samples reduced modulo $q$, using the convention of \autoref{subsec:lwe-dec}. A correctness inequality is a condition on these lifted coefficients, not a deterministic guarantee for an unbounded Gaussian distribution.

Suppose we have an $(n-1)$-degree polynomial $M \in \mathcal{R}_{\langle n, t \rangle}$ whose coefficients represent the plaintext numbers to encrypt.


\begin{enumerate}
\item Randomly pick a list of $k$ $(n-1)$-degree polynomials $\{A_i\}_{i=0}^{k-1} \xleftarrow{\$} \mathcal{R}_{\langle n,q \rangle}^{k}$ as a one-time public mask.
\item Randomly pick a small polynomial $E \xleftarrow{\chi_\sigma} \mathcal{R}_{\langle n,q \rangle}$ as a one-time noise, whose $n$ coefficients are small numbers in $\mathbb{Z}_q$ randomly sampled from the Gaussian distribution $\chi_\sigma$. 
\item Scale $M$ by $\Delta$, which is to compute $\Delta \cdot M$. This converts $M \in \mathcal{R}_{\langle n, t \rangle}$ into $\Delta \cdot M \in \mathcal{R}_{\langle n,q \rangle}$.
\item Compute $B = \sum\limits_{i=0}^{k-1}{(A_i \cdot S_i)} + \Delta \cdot M + E \in \mathcal{R}_{\langle n,q \rangle}$. 
\end{enumerate}

$ $

The GLWE encryption formula is summarized as follows:

$ $

\begin{tcolorbox}[title={\textbf{\tboxlabel{\ref*{subsec:glwe-enc}} GLWE Encryption}}]

\textbf{\underline{Initial Setup}:} $\Delta = \left\lfloor\dfrac{q}{t}\right\rfloor$, $\{S_i\}_{i=0}^{k-1} \xleftarrow{\$} \mathcal{R}_{\langle n, \textit{tern} \rangle}^k$

$ $

$ $

\textbf{\underline{Encryption Input}:} $M \in \mathcal{R}_{\langle n, t \rangle}$, $\{A_i\}_{i=0}^{k-1} \xleftarrow{\$} \mathcal{R}_{\langle n,q \rangle}^{k}$, $E \xleftarrow{\chi_\sigma} \mathcal{R}_{\langle n,q \rangle}$

$ $

\begin{enumerate}

\item Scale up $M \longrightarrow \Delta M \text { } \in \mathcal{R}_{\langle n, q\rangle}$

\item Compute $B = \sum\limits_{i=0}^{k-1}{(A_i \cdot S_i)} + \Delta  M + E \text{ } \in \mathcal{R}_{\langle n,q \rangle}$

\item $\textsf{GLWE}_{S,\sigma}(\Delta M + E) = (\{A_i\}_{i=0}^{k-1}, B) \text{ } \in \mathcal{R}_{\langle n,q \rangle}^{k + 1}$ 



\end{enumerate}

\end{tcolorbox}

\subsection{Decryption}
\label{subsec:glwe-dec}

\begin{enumerate}
\item Given the ciphertext $(\{A_i\}_{i=0}^{k-1}, B)$ where $B = \sum\limits_{i=0}^{k-1}{(A_i \cdot S_i)} + \Delta \cdot M + E \in \mathcal{R}_{\langle n,q \rangle}$, compute $B - \sum\limits_{i=0}^{k-1}{(A_i \cdot S_i)} = \Delta \cdot M + E$. 
\item Round each coefficient of the polynomial $\Delta \cdot M + E \in \mathcal{R}_{\langle n,q \rangle}$ to the nearest multiple of $\Delta$ (i.e., round it as a base $\Delta$ number), which is denoted as $\lceil \Delta \cdot M + E \rfloor_{\Delta}$. This operation successfully eliminates $E$ and gives $\Delta \cdot M$. One caveat is that $E$'s each coefficient $e_i$ has to be $|e_i| < \dfrac{\Delta}{2}$ to be eliminated during the rounding. Otherwise, some of $e_i$'s higher bits will overlap the plaintext $m_i$ coefficient's lower bit and won't be eliminated during decryption, corrupting the plaintext $m_i$.
\item Compute $\dfrac{\Delta \cdot M} {\Delta}$, which is equivalent to scaling down each polynomial coefficient in $\Delta \cdot M$ by $\Delta$. 
\end{enumerate}

$ $

In summary, the GLWE decryption formula is summarized as follows:

$ $

\begin{tcolorbox}[title={\textbf{\tboxlabel{\ref*{subsec:glwe-dec}} GLWE Decryption}}]
\textbf{\underline{Decryption Input}:} $\textsf{ct} = (\{A_i\}_{i=0}^{k-1}, B) \text{ } \in \mathcal{R}_{\langle n,q \rangle}^{k + 1}$

\begin{enumerate}
\item $\textsf{GLWE}^{-1}_{S,\sigma}(\textsf{ct}) = B - \sum\limits_{i=0}^{k-1}{(A_i \cdot S_i)} = \Delta  M + E \text{ } \in \mathcal{R}_{\langle n,q \rangle}$

\item Scale down 
$\Bigl\lceil \dfrac{ \Delta  M + E }{\Delta}\Bigr\rfloor \bmod t = M \text{ } \in \mathcal{R}_{\langle n, t \rangle}$

\end{enumerate}

For correct decryption, every noise coefficient $e_i$ of polynomial $E$ should be: $|e_i| < \dfrac{\Delta}{2}$.

\end{tcolorbox}


\subsubsection{Discussion}

The following size comparisons describe representation and brute-force search space only. Concrete security depends jointly on $n,k,q$, the error and secret distributions, and sample count; increasing one dimension is not by itself a security reduction.
\begin{enumerate}
\item \textbf{LWE} is a special case of GLWE where the polynomial ring's degree $n = 1$. That is, all polynomials in $\{A_i\}_{i=0}^{k-1}, \{S_i\}_{i=0}^{k-1}, E$, and $M$ are zero-degree polynomial constants. Instead, there are $k$ such constants for $A_i$ and $S_i$, so each of them forms a vector.
\item \textbf{RLWE} is a special case of GLWE where $k = 1$. That is, the secret key $S$ is a single polynomial $S_0$, and each encryption is processed by only a single polynomial $A_0$ as a public mask.
\item \textbf{Size of $\bm{n}$:} The polynomial $S_i = s_{i,0} + s_{i,1}X + \cdots + s_{i, n-1}X^{n-1} $ has $n$ coefficients. With independent uniform ternary coefficients, the complete secret has $3^{nk}$ possible values. This is a count of exhaustive-search candidates, not a lower bound on the best attack or a proof of the search assumption (\autoref{subsec:lattice-overview}). Higher-degree plaintext polynomials can represent more coefficients; the resulting throughput also depends on the operation costs.
%\item \textbf{Size of $\bm{q}$:} Increasing $q$ increases the ciphertext molulus space, making it hard to attack the cryptosystem. 
%Also, a large modulo space creates an enough separation between the plaintext bits and the noise bits in the ciphertext, which reduces the growth rate of the noise bits across ciphertext multiplications (will be discussed in \autoref{subsubsec:glwe-mult-plain-discussion}) and thus we can do a more number of ciphertext multiplications before the noise bits overflow to the plaintext bit area. 
\item \textbf{Size of $\bm{k}$:} Increasing $k$ adds secret polynomials $(S_0, S_1, \ldots, S_{k-1})$ and public mask polynomials $(A_0, A_1,\ldots, A_{k-1})$. Each ciphertext still has one message polynomial $M$ and one error polynomial $E$. Comparing two choices of $n,k$ requires their actual ring/module assumptions and distributions, not only the number of secret coefficients.
\item \textbf{Reducing the Ciphertext Size:} A seeded ciphertext can store a fresh independent public seed $d$ and $B$. A specified public expansion $G$ produces all $kn$ mask coefficients, with parameter/index domain separation and an unbiased coefficient mapping. Its input is never the secret. Because the seed is public, ordinary hidden-seed PRG security does not justify replacing $(d,G(d))$ by $(d,A)$ for independent uniform $A$: testing $A=G(d)$ distinguishes those distributions. For a fixed expansion, the necessary joint assumption concerns $(d,\langle G(d),S\rangle+E)$ versus $(d,V)$ for uniform $V$, with the actual secret/error distributions and sample count.
\end{enumerate}

\noindent\textbf{A precise seeded model:} If expansion is an ideal random oracle $H:\{0,1\}^h\to\mathcal R_{\langle n,q\rangle}^k$, a classical adversary making $Q$ encryption queries and $Q_H$ oracle queries has guessing bias
\[
\epsilon\le\delta(T',Q)+\min\left\{1,\frac{QQ_H+Q(Q-1)/2}{2^h}\right\}.
\]
Here $\delta$ is the actual-distribution decisional GLWE advantage, with the convention of \autoref{subsec:lattice-scheme}. To simulate a ciphertext, choose its seed independently and program the previously unassigned $H(d)$ with the uniform mask from a challenge sample; add the selected message to its body. The simulation fails only if the seed was queried earlier or repeats an earlier encryption seed. At the $j$th encryption at most $Q_H+j-1$ points are occupied, giving the displayed union bound. Outside that event real samples give real encryption; uniform sample bodies hide the message. Abort with an independent fair guess on failure. This proves the stated classical random-oracle result. It neither replaces a concrete XOF by an ideal oracle without an assumption nor proves the same bound for quantum oracle queries. See \href{https://pq-crystals.org/kyber/data/kyber-specification-round3-20210804.pdf}{the Kyber v3.02 specification, Sections 4.3.1 and 5.4}, for the corresponding distinction between public XOF expansion and secret pseudorandom generation. Reusing a mask also exposes the difference of two message phases plus their error difference, so seed freshness is essential.

\subsection{An Alternative Version of GLWE}
\label{subsec:glwe-alternative}

The following is an alternative version of \tboxlabel{\ref*{subsec:glwe-enc}}, where the sign of each $A_iS_i$ is flipped in the encryption and decryption formula as follows: 

\begin{tcolorbox}[title={\textbf{\tboxlabel{\ref*{subsec:glwe-alternative}} An Alternative GLWE Cryptosystem}}]

\textbf{\underline{Initial Setup}:} $\Delta = \left\lfloor\dfrac{q}{t}\right\rfloor$, $\{S_i\}_{i=0}^{k-1} \xleftarrow{\$} \mathcal{R}_{\langle n, \textit{tern} \rangle}^k$

$ $

\par\noindent\rule{\textwidth}{0.4pt}


\textbf{\underline{Encryption Input}:} $M \in \mathcal{R}_{\langle n, t \rangle}$, $\{A_i\}_{i=0}^{k-1} \xleftarrow{\$} \mathcal{R}_{\langle n,q \rangle}^{k}$, $E \xleftarrow{\chi_\sigma} \mathcal{R}_{\langle n,q \rangle}$

$ $

\begin{enumerate}
\item Scale up $M \longrightarrow \Delta M \text { } \in \mathcal{R}_{\langle n, q\rangle}$

\item Compute $B = \textcolor{red}{-\sum\limits_{i=0}^{k-1}{(A_i \cdot S_i)}} + \Delta  M + E \text{ } \in \mathcal{R}_{\langle n,q \rangle}$

\item $\textsf{GLWE}_{S,\sigma}(\Delta M + E) = (\{A_i\}_{i=0}^{k-1}, B) \text{ } \in \mathcal{R}_{\langle n,q \rangle}^{k + 1}$ 

\end{enumerate}

\par\noindent\rule{\textwidth}{0.4pt}


\textbf{\underline{Decryption Input}:} $\textsf{ct} = (\{A_i\}_{i=0}^{k-1}, B) \text{ } \in \mathcal{R}_{\langle n,q \rangle}^{k + 1}$

\begin{enumerate}
\item $\textsf{GLWE}^{-1}_{S,\sigma}(\textsf{ct}) = B \text{ } \textcolor{red}{ + \sum\limits_{i=0}^{k-1}{(A_i \cdot S_i)}} = \Delta  M + E \text{ } \in \mathcal{R}_{\langle n,q \rangle}$

\item Scale down 
$\Bigl\lceil \dfrac{ \Delta  M + E }{\Delta}\Bigr\rfloor \bmod t = M  \text{ } \in \mathcal{R}_{\langle n, t \rangle}$

\end{enumerate}

For correct decryption, every noise coefficient $e_i$ of polynomial $E$ should be: $|e_i| < \dfrac{\Delta}{2}$.

\end{tcolorbox}

Even if the $A_iS_i$ terms flip their signs, the decryption stage cancels out those terms by adding their equivalent double-sign-flipped terms; thus, the same correctness of decryption is preserved as in the original version.  





\subsection{Public Key Encryption}
\label{subsec:glwe-public-key-enc}

\textbf{Scope:} This subsection is restricted to $k=1$, the standard RLWE public-key structure, with sufficiently large power-of-two ring dimension $n$ and a parameter/distribution family for which the required decisional RLWE assumptions are justified. The vector notation below is retained with a single component. It is not a generic public-key construction for arbitrary GLWE dimensions or the $n=1$ LWE special case. Such generalizations require a suitable public matrix construction and their own security argument. As in the setup above, assume $t\mid q$ for the displayed decryption bound.


The encryption scheme in \autoref{subsec:glwe-enc} assumes that it is the secret key owner who encrypts each plaintext. In this section, we explain a public key encryption scheme in which we create a public key counterpart of the secret key. Anyone who knows the public key can encrypt the plaintext in such a way that only the secret key owner can decrypt it. The high-level idea is that a portion of the components to be used in the encryption stage is pre-computed at the setup stage and published as a public key. At the actual encryption stage, the public key is multiplied by an additional randomness ($U$) and added to additional noise ($E_1, \vec{E}_2$) to create unpredictable randomness for each encrypted ciphertext. The actual scheme is as follows: 

\begin{tcolorbox}[title={\textbf{\tboxlabel{\ref*{subsec:glwe-public-key-enc}} RLWE Public Key Encryption ($k=1$)}}]

\textbf{\underline{Initial Setup}:} 
\begin{itemize}
\item The scaling factor $\Delta = \left\lfloor\dfrac{q}{t}\right\rfloor$
\item The secret key $\vec{S} = \{S_i\}_{i=0}^{k-1} \xleftarrow{\$} \mathcal{R}_{\langle n, \textit{tern} \rangle}^k$
\item The public key pair $(\mathit{PK}_1, \vv{\mathit{PK}}_2) \in \mathcal{R}_{\langle n, q \rangle}^{k+1}$ is generated as follows:

$\vec{A} = \{A_i\}_{i=0}^{k-1} \xleftarrow{\$} \mathcal{R}_{\langle n, q \rangle}^k$, \text{ } $E \xleftarrow{\chi_{\sigma}} \mathcal{R}_{\langle n, q \rangle}$

$\mathit{PK}_1 = \vec{A} \cdot \vec{S} + E \in \mathcal{R}_{\langle n, q \rangle}$

$\vv{\mathit{PK}}_2 = \vec{A} \in \mathcal{R}_{\langle n, q \rangle}^{k}$

\end{itemize}

\par\noindent\rule{\textwidth}{0.4pt}

\textbf{\underline{Encryption Input}:} $M \in \mathcal{R}_{\langle n, t \rangle}$, \text{ } $U \xleftarrow{\$} \mathcal{R}_{\langle n,\textit{tern} \rangle}, \text{ } E_1 \xleftarrow{\chi_{\sigma}} \mathcal{R}_{\langle n,q \rangle}, \text{ } \vec{E}_2 \xleftarrow{\chi_{\sigma}} \mathcal{R}_{\langle n,q \rangle}^k$

$ $

\begin{enumerate}

\item Scale up $M \longrightarrow \Delta M \text { } \in \mathcal{R}_{\langle n, q\rangle}$

\item Compute the following: 

$B = \mathit{PK}_1\cdot U + \Delta  M + E_1 \text{ } \in \mathcal{R}_{\langle n,q \rangle}$

$\vec{D} = \vv{{\mathit{PK}}}_2 \cdot U + \vec{E}_2 \in \mathcal{R}_{\langle n,q \rangle}^{k}$ \textcolor{red}{ $\rhd$ $\vv{\mathit{PK}}_2 \cdot U$ multiplies each element of $\vv{\mathit{PK}}_2$ by $U$}

\item $\textsf{GLWE}_{S,\sigma}(\Delta M + E_{\mathit{all}}) = (\vec{D}, B) \text{ } \in \mathcal{R}_{\langle n,q \rangle}^{k+1}$  \textcolor{red}{ $\rhd$ where $E_{\mathit{all}} = E \cdot U + E_1 - \vec{E}_2\cdot\vec{S}$}



\end{enumerate}


\par\noindent\rule{\textwidth}{0.4pt}

\textbf{\underline{Decryption Input}:} $\textsf{ct} = (\vec{D}, B) \text{ } \in \mathcal{R}_{\langle n,q \rangle}^{k + 1}$

\begin{enumerate}
\item $\textsf{GLWE}^{-1}_{S,\sigma}(\textsf{ct}) = B - \vec{D} \cdot \vec{S} = \Delta  M + E_{\mathit{all}} \text{ } \in \mathcal{R}_{\langle n,q \rangle}$

\item Scale down 
$\Bigl\lceil \dfrac{ \Delta  M + E_{\mathit{all}} }{\Delta}\Bigr\rfloor \bmod t = M  \text{ } \in \mathcal{R}_{\langle n, t \rangle}$

\end{enumerate}

For correct decryption, every noise coefficient $e_i$ of polynomial $E_{\mathit{all}}$ should be: $|e_i| < \dfrac{\Delta}{2}$.

\end{tcolorbox}


The equation in the 1st step of the decryption process is derived as follows:

$\textsf{GLWE}^{-1}_{S,\sigma}\bm{(}\text{ } \textsf{ct} = (\vec{D}, B) \text{ }\bm{)} = B - \vec{D}\cdot\vec{S}$

$ = (\mathit{PK}_1\cdot U + \Delta  M + E_1) - (\vv{\mathit{PK}}_2 \cdot U + \vec{E}_2)\cdot \vec{S} $

$=  (\vec{A} \cdot \vec{S} + E)\cdot U + \Delta M + E_1 - (\vec{A}\cdot U)\cdot \vec{S} - \vec{E}_2\cdot\vec{S}$

$=  (U\cdot\vec{A}) \cdot \vec{S} + E \cdot U + \Delta M + E_1 - (U\cdot \vec{A}) \cdot \vec{S} - \vec{E}_2\cdot\vec{S}$

$= \Delta M + E \cdot U + E_1 - \vec{E}_2\cdot\vec{S}$

$= \Delta M + E_{\mathit{all}}$ \textcolor{red}{ $\rhd$ where $E_{\mathit{all}} = E \cdot U + E_1 - \vec{E}_2\cdot\vec{S}$} 


$ $

\para{Security:} For $k=1$, write the public key as $(A,P)$ with $P=AS+E$, and the ciphertext as $(D,B)=(AU+E_2,PU+E_1+\Delta M)$. Assume decisional RLWE for the actual distribution of $S$ and error $E$, and for the actual distribution of $U$ with two samples whose errors have the distributions of $E_2,E_1$. In this subsection both secrets are ternary and all three error laws are the same, so one corresponding two-sample RLWE assumption covers both uses. It does not follow merely from a theorem stated for Gaussian secrets.

The single-challenge IND-CPA proof has two hybrids. First replace $P=AS+E$ by an independent uniform ring element. A reduction embeds one RLWE sample as the public key and can perform every encryption itself by choosing fresh $U,E_1,E_2$; it does not need $S$. Next, with uniform $(A,P)$, replace the challenge's $(D,B-\Delta M)$ by independent uniform ring elements. A reduction takes two samples $(A,Z_2),(P,Z_1)$ sharing the same unknown secret $U$, uses their masks as the public key, and returns $(Z_2,Z_1+\Delta M)$ once the challenge messages are chosen. Real samples give $Z_2=AU+E_2$ and $Z_1=PU+E_1$; uniform samples hide the chosen message completely. Other encryptions are publicly computable with fresh known randomness. Thus, with the guessing-bias convention of \autoref{subsec:lattice-scheme},
\[
\epsilon_{\mathrm{PKE}}\le\delta_S(T'_1,1)+\delta_U(T'_2,2).
\]
The two terms use the stated actual distributions and include reduction overhead. This is a complete conditional CPA reduction for the scoped construction; it does not certify a numerical parameter family or secret-dependent public-key encryption queries. See \href{https://web.eecs.umich.edu/~cpeikert/pubs/ideal-lwe.pdf}{Lyubashevsky--Peikert--Regev, Section 1.1, Applications}. For an unrestricted vector version, the second hybrid would require rank-one RLWE with secret $U$ and $k+1$ samples, not merely the rank-$k$ GLWE assumption used in the first hybrid.

The cancellation calculation proves a separate correctness identity with $E_{\mathit{all}}=EU+E_1-E_2S$.

For lifted coefficients in the negacyclic ring,
\[
\|E_{\mathit{all}}\|_\infty\le\|E\|_\infty\|U\|_1+\|E_1\|_\infty+\|E_2\|_\infty\|S\|_1.
\]
The right-hand side being less than $\Delta/2$ suffices for decryption when $t\mid q$.

The unrestricted construction would be insecure for $n=1$ and ternary $U$, even if the LWE vector dimension $k$ were large: an attacker need only try $U\in\{-1,0,1\}$ and test whether all $D_i-A_iU$ are small, then decode $B-\mathit{PK}_1U$. In particular, $U=0$ occurs with probability $1/3$ and leaves $B=\Delta m+E_1$. This is why the general GLWE security claim cannot be inferred from the displayed algebra.
